\documentclass[11pt, oneside]{article} 
\usepackage{geometry}
\geometry{letterpaper} 
\usepackage{graphicx}
	
\usepackage{amssymb}
\usepackage{amsmath}
\usepackage{parskip}
\usepackage{color}
\usepackage{hyperref}

\graphicspath{{/Users/telliott/Dropbox/Github-Math/figures/}}
% \begin{center} \includegraphics [scale=0.4] {gauss3.png} \end{center}

\title{Quadratics, gently}
\date{}

\begin{document}
\maketitle
\Large

%[my-super-duper-separator]

\section*{Summary}

This section summarizes what you want to remember after you're finished reading.  Sometimes called TL;DR (too long didn't read).  Skim or skip this part if you are new to the topic.

The general equation for a quadratic is
\[ y = ax^2 + bx + c \]
where $a$, $b$ and $c$ are constants. 

$\circ$ \ $a$ determines how steeply the sides of parabola curve and its sign tells in what direction the "cup opens", up or down.

$\circ$ \ $c$ shifts the parabola's vertical position and is also the $y$-value when $x = 0$.

$\circ$ \ $b$ and $a$ together determine where the pointy end lies.  

This point is called the vertex of the parabola and it is always the minimum for $y$ (or maximum if the parabola opens down).  We will call its $x$-value $m$ for minimum, and just give the result, for now.
\[ m = - \frac{b}{2a} \]

To "solve" a quadratic means to find its zeros, or roots.  These are the $x$-values that make the whole thing equal to zero.
\[ ax^2 + bx + c = 0 \]
Dividing by $a$ does not change the roots.  The same values of $x$ satisfy
\[ x^2 + \frac{b}{a} x + \frac{c}{a} =  0 \]
The zeros can always be found using the quadratic formula.  Here is a simple version of that:
\[ x = m \pm \ \sqrt{m^2 - \frac{c}{a}} \]

You must memorize both the formula for $m$ and this one.

Sometimes we can guess the zeros since
\[ (x - s)(x - t) = x^2 - (s + t)x + st = 0 \]
See if you can two factors that add to give negative $b$ and multiply to give $c$.

As an example
\[ x^2 + 2x - 63 = y \]
Find integers two units apart that multiply to give $63$, then play with the signs.  The roots are $-7$ and $9$.  The vast majority of equations do not have integer roots, but this exercise does focus attention on what they are.

The easiest derivation of the quadratic formula is by "completing the square".  Start with
\[ x^2 + \frac{b}{a} x + \frac{c}{a} =  0 \]

We will first simplify the middle term, the cofactor of $x$.  Remember that $m = -b/2a$ so $-2m = b/a$ and then plug in to get
\[ x^2 - 2mx + \frac{c}{a} =  0 \]

The bright idea is to turn this into $(x - m)^2$ which we do by adding and subtracting $m^2$:
\[ x^2 - 2mx + m^2 - m^2 + \frac{c}{a} = 0 \]
\[ (x - m)^2 - m^2 + \frac{c}{a} = 0 \]
The rest is simple algebra, which I leave to you.

\subsection*{Introductory examples}

A quadratic equation contains the square of a variable multiplied by some constant, which may be $1$.  Often we call the variables $y$ and $x$, so we have something like
\[ y = ax^2 \]

It may be that $a = 1$, and then $y = x^2$, but $a$ cannot be zero, or we don't have a quadratic.  

Then in addition there may be an optional term containing the first power of $x$ and also perhaps a final constant.  The general form is
\[ y = ax^2 + bx + c \]
where $a, b$ and $c$ are constants.  $b$ or $c$ (or both) might be zero.

The two most common applications are maximization problems, and gravity.

\subsection*{maximizing area}
We need to build a fence or a wall from a fixed amount of material.  There are constraints:   it must be in the shape of a rectangle, and we want the area enclosed to be a maximum.

Although you will usually be given some definite amount of material, like $100$ feet of fencing, but it does no harm to re-scale our problem, as long as we understand how to undo the scaling once we have the answer.

Let the perimeter of the rectangle be $2$, the half-perimeter is $1$, and then if one dimension, say the short side of the rectangle is $h$, the width is $1 - h$, so the area is
\[ A = hw = h(1 - h) = -h^2 + h \]

As will be obvious later, since $a < 0$, the graph of this equation forms a "cup" that opens down.  The value of $A$ at the vertex is a maximum.

Two fundamental facts about parabolas:   

$\circ$ \ the maximum (or minimum) always occurs at the vertex

$\circ$ \ the $x$-coordinate of the vertex has the value $-b/2a$

So the simplest parabola $y = x^2$ has $b = 0$, because there is no term like $bx$, and the vertex is at $x = 0$, $y = 0$.

In our problem $a = -1$ and $b = 1$, the vertex is the maximum and it is at $h = 1/2$.   Since $h + w = 1$, it follows that $h = w = 1/2$.

The meaning of this answer is that the square ($h = w$) has the maximum area for a rectangular shape with a fixed perimeter.

Actually, you may know that if the shape is not rectangular the maximum area is for a circle or disk, but that is not the problem we're solving!

\subsection*{gravity}

A second application is for objects falling under the influence of gravity.  Gravity is an acceleration, and that acceleration is constant, which means that if the speed or velocity after falling for one second is $v$, the velocity after two seconds is twice $v$ and so on.  

So the distance $d = vt$ covered from zero time up to 1 second is less than that covered between $t = 1$ and $t = 2$ seconds.  

Let me just give you the equation.  If we drop our rock from the Tower of Pisa and the tower stands a height $h_0$ above the ground, the distance $h$ above the ground after $t$ seconds is
\[ h = -16t^2 + h_0 \]
This makes sense.  At $t = 0$, $h = h_0$ and then $h$ will be smaller than $h_0$ by some amount $-16t^2$ when $t > 0$.

If the initial height $h_0$ is $144$ feet, then a dropped rock will hit the ground at $h=0$ which means that
\[ 0 = -16t^2 + 144 \]
\[ t = \sqrt{144/16} = \sqrt{9} = 3 \text{ sec} \]

We take only the positive square root, but to explain that will take more time than we want for this example right now.  We'll come back to it.

There is a limit to the final, or terminal velocity.  It comes not from gravity but from air resistance.  Someone who jumps from an airplane or gets thrown out through a window of the Nakatomi Plaza has a terminal velocity of about $180$ feet per second or $120$ miles per hour.

\section*{Basic theory}

All that is to justify why we want to have an introduction to quadratic equations.  The first part of what's coming is a really gentle start, the second section gives some more detail about the theory, and the third works through some more examples.  

\subsection*{introduction}

You know that a linear equation is something like $y = mx + b$, where $x$ and $y$ are the variables, and $m$ is the slope of the line that is the graph of the equation. 

When  $x = 0$, then $y = b$.  

$b$ is the value that $y$ has when the line crosses the $y$-axis, where $x = 0$, called the $y$-intercept.  It's a reasonable idea to use the symbol $y_0$ instead of $b$, since we'll be using $b$ in the quadratic.  Say "y-naught" as in the $y$ when $x$ is naught, or zero.

Since we are also going to use the symbol $m$ for something else today, allow me to write $y = kx + y_0$, with $k$ being the slope.

We see that a linear equation contains $kx$ or $kx^1$ ($x$ to the power $1$), and it may have a constant as well, which you can think of as the cofactor for the $x^0$ power.

An extension to higher powers is natural.  We write
\[ y = ax^2 + bx + c \]

where $x$ and $y$ are again variables and $a$, $b$ and $c$ are constants.  I always think of this as the \emph{standard form} of a quadratic equation, although now that I look it up using Google it says that standard form is
\[ 0 = ax^2 + bx + c \]

and one page even says something else.  That's the problem with the internet --- someone is always wrong.

Anyway, for us here, standard form means that we write out all three cofactors, $y = ax^2 + bx + c$.

Here's an example
\[ y = x^2 + 2x - 3 \]
an equation whose graph looks like this in Desmos:

\begin{center} \includegraphics [scale=0.5] {quad1.png} \end{center}

We might notice a couple of things about this graph.  First, it is not a straight line as with the linear equation, but a curve that bends upward.  The graph is said to resemble a "cup" that "opens upward".  If the leading constant $a$ is negative then the graph will flip directions, and the "cup" opens downward.

Second, this curve crosses the $x$-axis, which means that for some values of $x$, the result of the equation is $y = 0$.

There is a minimum value for $y$ on the curve, which occurs at a particular value of $x$, we will call that value $x = m$.

Finding $x$ which results in the minimum is the "hard" part, though it is not exactly hard.  If you need the actual value of $y$ at the minimum, plug $m$ into the quadratic equation.

The point on the curve where $x=m$ is called the \emph{vertex}.  If the cup would open down instead ($a < 0$), then this value would be a maximum.

A very important fact is that the zeros of the equation (the $x$-values where $y = 0$) are evenly spaced on either side of $m$, so we can think of $m$ standing for mean, the average of the two zeros.

In fact, the curve as a whole is symmetric about the vertical line $x = m$.  That line is called the \emph{axis of symmetry}.

It is worth emphasizing that if we know the zeros $s$ and $t$ then we know $m$ because it always lies exactly halfway in between.
\[ m = \frac{s + t}{2} \]

\subsection*{simplification}

At various points in what comes next, we may modify the general equation 
\[ y = ax^2 + bx + c \]
to make things simpler by setting $a = 1$
\[ y = x^2 + bx + c \]
or for the specific case when $y = 0$
\[ 0 = x^2 + \frac{b}{a} x + \frac{c}{a} \]
or even go further sometimes by setting $b = 1$
\[ y = x^2 + x + c \]
and then, having introduced an idea, go back to considering what happens in the general case.

\subsection*{examples}
Here are some simple equations to plug into Desmos and look at how the graph changes.

\[ y = x^2 \]
\[ y = ax^2 \]

Add a slider for $a$ and check out what happens when you change it.  $a$ is called the shape factor.  If $b$ is not zero, then the role of $a$ is a little more complicated because the position of the vertex is proportional to $b/a$.  When $b = 0$ as for $y = ax^2$, this doesn't happen.

In this case, the graph does not really cross the $x$-axis but just touches it at the vertex.  The vertex lies on the $x$-axis and the $y$-value corresponding to $x = 0$ is $y = 0$.

Next enter the equation
\[ y = x^2 - 2x + c \]
Add a slider for $c$ and notice what happens when you change it.   

Adding $c$ simply moves the curve up or down vertically without changing its shape or the axis of symmetry $x = m$.  It also changes the $x$-intercept.
\begin{center} \includegraphics [scale=0.5] {quad2.png} \end{center}

When $c=-3$
\[ y = x^2 + 2x - 3 \]

The zeros (the values of $x$ giving $y = 0$ are $x = -3, +1$.  We can actually read them off the graph.  You can check the arithmetic.  $(-3)^2 - 6 - 3 = 0$.  Also, $1^2 + 2 - 3 = 0$.

For any quadratic, there are some values of $c$ (more positive) where the curve won't cross the $x$-axis at all.  This means that there simply aren't any values of $x$ that give $y = 0$.

Here the vertex is $4$ units down from the $x$-axis so if we were to add $5$ to what we had before, giving
\[ y = x^2 + 2x + 2 \]
If you check in Desmos you'll find that this curve does not cross the $x$-axis.

And if we were to add only $4$ we would have
\[ y = x^2 + 2x + 1 = (x + 1)^2 \]
Each of these curves has its minimum value at $x = m = -1$.  For this particular example the minimum $y = (-1 + 1)^2 = 0$.  The parabola just touches the $x$-axis at the point $(-1,0)$.

The role of $b$ is a bit tricky to analyze, but you can take a look in Desmos with a slider.  Here's a challenge:  what shape do you think the vertex traces out when $b$ varies?

\subsection*{second form of the equation}
Let us introduce another way of writing a quadratic and give a specific example:
\[ y = (x - s)(x - t) \ \ \ \ \ \ \ \ \ \ y = (x + 3)(x - 1) \]

$s$ and $t$ are numbers, but not just any numbers.  When $x = s$ or $x = t$, then $y = 0$.

Multiply out the equation on the right to obtain 
\[ y = x^2 + 2x - 3 \]

You can see why the graph of this equation, which we looked at before, crosses the $x$-axis at $-3$ and $1$.  When $x = -3$ or $x = 1$ is plugged into $y = (x + 3)(x - 1)$, we quickly get $y = 0$, without doing any real arithmetic.

A neat fact about the vertex is that $m$, the $x$-value of the vertex, always lies halfway on the $x$-axis between $s$ and $t$.  $m$, the $x$-coordinate of the vertex, is the \emph{average} of $s$ and $t$.
\[ m = \frac{s + t}{2} = \frac{-3 + 1}{2}  = -1 \]

You might be worried about the situation when the graph never crosses the $x$-axis. In that case, what could $s$ and $t$ be to give $y = 0$?  Does $y = (x - s)(x - t)$ even make sense in that case?  

It turns out to have a deep meaning in terms of what are called complex numbers, involving $\sqrt{-1}$, but I think it's better to postpone thinking about it for now.  

When talking about roots, we will agree to only worry about functions whose graphs actually cross the $x$-axis or at least touch it (like $y =x^2$).

And by the way, for $y = x^2$ we can say that it has duplicate roots, both $x = 0$, since $(x - 0)(x - 0)$ is the equation written in the $s,t$ form.

\subsection*{a third form}
There is  yet a third common way to write the equation of a parabola.  We introduce it because it will justify the following fact, which you must memorize:  for any parabola
\[ y = ax^2 + bx + c \]
the $x$-coordinate of the vertex is at
\[ m = -\frac{b}{2a} \]

If a problem asks for the minimum or maximum only,  and you have the equation, you can answer without any work.

Now, suppose the graph does not have its vertex at the origin but instead at some other point like $(3,-2)$.  Is there a formula for that?  Look at this.
\begin{center} \includegraphics [scale=0.5] {quad3.png} \end{center}

We used Desmos to verify that the graph of the equation $y + 2 = (x - 3)^2$ has its vertex at $(x = 3, y = -2)$.  This is not a coincidence!  Try it with any other pairs of numbers, even decimals or fractions.

It is traditional to use the notation $x = h$ and $y = k$ for the coordinates of the vertex.  In general, I prefer $x = m$, for the reasons given before.  However in this section we use $h$ and $k$.

When the vertex is at $(h,k)$, the equation is $y - k = (x - h)^2$.  Pay attention to the minus signs.
\[ y - 1 = (x - 2)^2 \ \ \ \ \ \rightarrow \ \ \ \ \text{vertex } (2,1) \]
\[ y + 6 = (x - 3)^2 \ \ \ \ \ \rightarrow \ \ \ \ \text{vertex } (3,-6) \]
\[ y = (x + 3)^2 \ \ \ \ \ \rightarrow \ \ \ \ \text{vertex } (-3,0) \]

You can just memorize this fact and use it, but if you want to have a way to justify it, think of using $h$ and $k$ to adjust the vertex to the origin $(0,0)$.  If the vertex was originally at $x = 3$, we subtract $3$ from \emph{every} $x$ and that moves the vertex to $x = 0$.

If we take all the $(x,y)$ values on the parabola and shift them by subtracting $(3,-2)$, the curve is simply moved or \emph{translated} to have its vertex at the origin without changing the shape.

Let us first think about the case where the shape factor $a = 1$.  Then
\[ y - k = (x - h)^2 = x^2 - 2xh + h^2 \]
Compare that to the standard form (with $a=1$):
\[ y = x^2 + bx + c \]

We have written the same equation in two different forms.  The cofactors must match!
\[ bx = - 2hx  \ \ \ \ \ \  \ \ \ \ \  b = - 2h \]

Rearranging, $h = -b/2$.  That's pretty simple.

If there is a shape factor $a$ it just multiplies the whole thing as in $y - k = a(x - h)^2$.  Again, we can play with the equation by multiplying out
\[ y - k = a(x - h)^2 \]
\[ = ax^2 - 2axh + ah^2 \]

If you compare with the general form of the equation
\[ y = ax^2 + bx + c \]

Again, the cofactors (including for $x$) must match.  That is
\[ bx = - 2ahx \]
\[ h = - \frac{b}{2a} \]

This says that for \emph{any} quadratic, the $x$-coordinate of the vertex, which we also called $m$ before, can be found from the general form

\[ ax^2 + bx + c  \ \ \ \ \ \ \Rightarrow  \ \ \ \ \ m = -b/2a \]

This is fundamental and definitely worth memorizing.  The most basic thing we are asked to do with a quadratic is to find the position of the vertex, which always gives the minimum (or maximum) value for $y$.

If you also want the $y$-value, just plug $x = -b/2a$ into the equation and do the arithmetic.

\subsection*{zeros}
The second fundamental values we need to obtain for a quadratic are the zeros, $x$ values which make $y = 0$ (provided they exist).

We saw a good example before, with
\[ y = (x + 3)(x - 1) = x^2 + 2x - 3 \]
Then, the values of $x$ which give $y = 0$ were just $x = -3$ and $x = 1$.  

For the general case, we will need to go backward starting from $ax^2 + bx + c$ and find something like $(x - s)(x - t)$.  

In simple cases we might be able to look at $x^2 + bx + c$ and see the answer.

\[ x^2 + 3x + 2 = 0 \]
Notice that $1 + 2 = 3$ and $1 \cdot 2 = 2$.  So we are looking for two factors which add together give to $b$ and multiply together to give $c$.
\[ = (x + 1)(x + 2) \]

While you can get good at this, with practice, most equations cannot be solved by this method, because almost always the roots are not integers.  They might be fractions, or square roots or even something else, as we mentioned.

Here are a few more examples of roots with the resulting $b,c$ values:
\[ (x + 1)(x - 5) \Rightarrow b = -4, c = -5 \]
\[ (x - \frac{1}{2})(x - \frac{1}{4})) \Rightarrow b = - \frac{3}{4}, c = \frac{1}{8} \]

Again, the standard form with $y = 0$ is
\[ ax^2 + bx + c = 0 \]

An important fact is that factoring out $a$
\[ a(x^2 + \frac{b}{a} x + \frac{c}{a}) = 0 \]
or even multiplying both sides by $1/a$
\[ x^2 + \frac{b}{a} x + \frac{c}{a} = 0 \]
\emph{doesn't change the roots}.

You should play with an example in Desmos to see this.  Changing $a$ will change the shape but not the roots.

This is like writing
\[ y = a(x - s)(x - t) \]

$a$ is a non-zero constant, so the only way the above equation can be zero is if $x = s$ or $x = t$ and this equation has the same roots as the original one.

Here's an example from a problem we will see later, we need to find the roots of
\[ -16t^2 + 64t + 96 = 0 \]
Because of that $0$ on the right-hand side, it is perfectly legal to factor out the leading term of $(-16)$ and then
\[ = (-16) (t^2 - 4t - 6) = 0 \]

The zeros of the first are the same as the zeros of the second
\[ t^2 - 4t - 6 = 0 \]

Unfortunately, we need the quadratic equation to actually solve this one.

\section*{Theory}

\subsection*{quadratic formula}

The general formula to find the roots of a quadratic equation is a somewhat complicated looking beast.
\[ x = \frac{-b \pm \ \sqrt{b^2 - 4ac}}{2a} \]

Breaking it into two separate terms we have
\[ x = - \frac{b}{2a} \pm \frac{\sqrt{b^2 - 4ac}}{2a} \]


You may recognize, however, that the first part is just the $x$-value of the vertex, $m = -b/2a$.

In fact, if you bring the denominator up under the square root
\[ x =  - \frac{b}{2a} \pm \sqrt{\frac{b^2}{4a^2} - \frac{c}{a}} \]
\[ = m \pm \sqrt{m^2 - \frac{c}{a}} \]

That looks like something I can remember.

A second approach that could help is to break it down a different way.  The part under the square root is called the discriminant $D$.
\[ D = b^2 - 4ac \]

Notice that if $b^2 < 4ac$ then $D = b^2 - 4ac < 0$ and \emph{and then} we have a problem with $\sqrt{D}$.  

This is what happens in the case where the graph of the equation does not cross the $x$-axis.  Remember when we said that making $c$ more positive eventually results in a graph with no roots.  Look at $D$ to see why.

Let's not worry about that now.  The main thing is that we can write the quadratic formula in two parts as $D = b^2 - 4ac$ and then
\[ x = \frac{- b \pm \ \sqrt{D}}{2a} \]

or even better
\[ x = m \pm \ \frac{\sqrt{D}}{2a} \]
Maybe that will help.  Otherwise, write it down, put it somewhere on your desk, and do a bunch of problems to help engrave it on the surface of your brain.

After years of practice, I still wake up from sleep reciting: "negative $b$, plus or minus the square root ... of $b^2 - 4ac$, all over 2$a$."

\subsection*{easier method}

At this point, the question arises of how to derive the quadratic formula which we gave above in the classic formulation and in a simpler version in the summary at the start of this write-up.  I like the simple version:
\[ x = m \pm \sqrt{m^2 - \frac{c}{a}} \]

I know two approaches to the derivation:  the first is "completing the square".   We have already done that, using a small trick (substitution of $m$ for $-b/2a$) to make it clearer.  

The second method is to analyze the zeros as being equidistant from $m$ and doing a comparison of forms such as we used in originally finding a value for $m$.

Remember that the problem we want to solve is
\[ ax^2 + bx + c  = 0 \]
that is, we want the values of $x$ that make this true.  As we said, the same $x$ values solve
\[ x^2 + \frac{b}{a} + \frac{c}{a} = 0 \]

Look at the symmetry of the graph.  For any particular value of $y$ draw the horizontal line to find the two places where $x$ gives that result.  For each pair of $x$-values, they will be symmetric to $m$, meaning that if the pair is $x_1, x_2$ then the distances are equal
\[ m-d = x_1 \ \ \ \ \ \ \ \ \ m + d = x_2 \]
\[ 2m = x_1 + x_2 \]

$m$ is the average of $x_1$ and $x_2$.

As a specific example of this, the roots, call them $s$ and $t$, are symmetric with $m$.
\begin{center} \includegraphics [scale=0.5] {quad2.png} \end{center}

Suppose that $s$ is smaller than $m$ and $t$ is the larger one.  Then, the distance between $s$ and $m$ or between $t$ and $m$ is $d$ and
\[ m - d = s \]
\[ m + d = t  \]
The change in sign happens because $m$ lies \emph{between} $s$ and $t$.  

So 
\[ (m - d)(m + d) = st \]
\[ m^2 - d^2 = st \]
\[ d^2 = m^2 - st \]

We want to know $d$.  We already know $m$, and we also know $st$.  

How is that?  Recall how we found $m = -b/2a$.  We compared
\[ x^2 - (s + t) x + st \]
\[ x^2 + (b/a) x + c/a \]
$m$ is the average of $s + t$ so $-2m = b/a$ and then $m = -b/2a$.

Now we focus on the fact that $c/a = st$.
\[ d^2 = m^2 - st = m^2 - c/a \]
\[ d = \sqrt{m^2 - c/a} \]

Then the roots are just $m \pm d$ or
\[ m \pm \sqrt{m^2 - \frac{c}{a}} \]
gives the zeros.

\subsection*{completing the square}
Here we derive the \emph{quadratic formula} by the standard approach, in all its glory.

We have a general equation
\[ ax^2 + bx + c = y \]
and we want to find the values of $x$ which give $y = 0$, the roots.
\[ ax^2 + bx + c = 0 \]
\[ x^2 + \frac{b}{a} x + \frac{c}{a} = 0 \]

Do this
\[ x^2 + \frac{b}{a}x + (\frac{b}{2a})^2 + \ [ \frac{c}{a} - (\frac{b}{2a})^2 \ ] \ = 0 \]
\[ (x + \frac{b}{2a})^2 + \ [ \frac{c}{a} - (\frac{b}{2a})^2 \ ] \ = 0 \]

Now it's just algebra.  If you recall that $m = -b/2a$ substitute into the above equation to obtain
\[ (x - m)^2 + \ [ \ \frac{c}{a} - m^2 \ ] \ = 0 \]

\[ (x - m)^2 = m^2 - c/a \]
\[ x - m = \pm \ \sqrt{m^2 - c/a} \]
\[ x = m \pm \ \sqrt{m^2 - c/a} \]

If you insist on the official formula, plug back in for $m = -b/2a$:
\[ x = -\frac{b}{2a} \pm \ \sqrt{(\frac{b}{2a})^2 - c/a} \]
\[ = -\frac{b}{2a} \pm \ \frac{1}{2a} \sqrt{b^2 - 4ac} \]

\section*{Examples}

\subsection*{maximizing profit}
You are the producer of a Broadway show.  The theater seats 120 people and you have priced the tickets at \$10 each.  You always fill the theater.

It costs you \$200 to rent the theater and another \$200 to pay the cast, so the fixed costs $F$ are $400$ and your profit for each show is
\[ P = nT - F \]
\[ = 120 \cdot 10 - 400 = 800 \]

\$800 per show.  That seems pretty good.

Your cousin Vinnie has an idea.  If you raise the price you'll make more money, but you'll also sell fewer tickets.  Acme Consulting has done a study which says that for every \$1 increase in the ticket price you'll have two empty seats.

As a red-blooded American, \emph{of course} you want the maximum profit.  What should you do?  There's clearly money to be made, if you raise the ticket price \$1, you'll sell 119 seats at \$11 for \$1309, giving a net profit of \$909.

Call the price increase $x$.  You will sell $n - 2x$ seats at a ticket price of $T + x$ for a profit
\[ P = (n - 2x)(T + x) - 400 \]
\[ = -2x^2 + (n - 2T)x + nT - 400 \]

The $x$ corresponding to the maximum of this quadratic is $-b/2a$.

\[ x  = -(n - 2T)/2(-2) \]
\[   = (120 - 20)/4 = 25 \]

You increase the ticket price from \$10 to \$35.  The profit will be
\[ P = (n - 50)(T + 25) - 400 \]
\[ = 70 \cdot 35 - 400 = 2050 \]
That's an increase of \$1250 per show.

\subsection*{odd numbers}
The product of two consecutive positive odd numbers is $255$.  Find the numbers.

Odd numbers can be described as $2n+1$ for some $n$, since that is one more than the even number $2n$.  So we have
\[ (2n + 1)(2n + 3) = 255 \]
\[ 4n^2 + 8n + 3 = 255 \]
\[ 4n^2 + 8n - 252 = 0 \]
We are looking for zeros.  So we can divide by $4$:
\[ n^2 + 2n - 63 = 0 \]
Two numbers multiply to give $-63$ and add to give $2$:
\[ (n + 9)(n - 7) = 0 \]
The roots are $-9, 7$.  But we specified positive numbers which means $n = 7$.  The numbers are $2 \cdot 7 + 1 = 15$ and $17$, easily confirmed.

Alternatively we can write the numbers as $2n + 1$ and $2n - 1$ and then
\[ (2n + 1)(2n - 1) = 255 \]
\[ 4n^2 - 1 = 255 \]
\[ n^2 = 64 \]
\[ n = \pm 8 \]
Again $n > 0$ so $n = 8$ and the numbers are $2n - 1 = 15$ and $2n + 1 = 17$.

\subsection*{gravity}
Many problems for quadratics involve gravity because the distance an object falls under gravity goes like $t^2$.  That was Galileo's famous discovery.  He used an inclined ramp and a water clock.  

The business about dropping objects of different weights off the Tower of Pisa was a popular activity during the century preceding Galileo.  His contribution was a thought experiment that is even more interesting.

\url{https://en.wikipedia.org/wiki/Galileo%27s_Leaning_Tower_of_Pisa_experiment}

In the metric system, the acceleration due to gravity (near the earth's surface) is very close to $10$ meters per second$^2$.  So, if we drop a rock, the distance traveled goes like 
\[ d = \frac{1}{2} 10 \ t^2 = 5 \ t^2 \]
After $1$ second, $5$ meters, after $2$ seconds $20$ meters, and so on.  If we know the distance to the ground and want the time:
\[ t = \sqrt{\frac{d}{5}} \]

Some problems involve throwing an object up in the air to begin with, then we need the general equation for movement in one dimension.
\[ x = \frac{1}{2}at^2 + v_0 t + x_0 \]

Which letter to use for the variable is your decision and really makes no difference.  We can use $y$ for height, by analogy with our graphs on the $xy$-axes.  We can use $h$, for height, or $d$ for distance.  Or we can just use $x$, as we might in outer space.

$x$ is really $x(t)$, the $x$-position depends on, it is \emph{a function of}, the time.  

The distance depends on the initial velocity $v_0$ times the time, which makes sense, and on the initial position $x_0$.

To use this equation we must keep our directions straight.  Gravity accelerates downward, normally initial heights will be positive, and initial velocity is usually positive but may be negative.  And of course the velocity may start positive but then surely turns negative, because gravity always wins.

You might be wondering about the other term, so consider a situation where $v_0$ and $x_0$ are both zero.  Then
\[ x = \frac{1}{2}at^2 \]

There's a bit of a mystery here, why do we have the factor of $1/2$, given the $v = at$?  

The $v$ in this formula is the velocity at the end of the interval, after $a$ has acted to boost the velocity from $v_0$ to $v$.

The correct relationship \emph{really} is
\[ x = \bar{v} t \]
where $\bar{v}$ is the average velocity
\[ \bar{v} = \frac{1}{2}at \]

We explore this a bit more in the next section.

\subsection*{time-dependence}
Distance equals velocity times time.

This is easy if the velocity is constant.  Travel west on the interstate at exactly $60$ miles per hour for $2$ hours and your distance will be $120$ miles from where you started (provided you don't start in Los Angeles).  

It is standard to use $s$ to refer to the distance traveled and $v$ for velocity.  If the velocity is constant then:
\[ s = vt \]
According to the internet, $s$ is from the Latin "spatium", for "space, room, or distance."

Suppose we plot velocity as a \emph{function of time} with $v$ on the $y$-axis and $t$ on the $x$-axis.

\begin{center} \includegraphics [scale=0.35] {velocity_time_3.png} \end{center}

Since the velocity is constant, the result is a straight horizontal line.  

Furthermore, the distance traveled is the \emph{area under the curve} (and above the $x$-axis) which is the area of a rectangle with sides $v$ and $t$ and as we said 
\[ s = vt \]

However, for most interesting problems the velocity is not constant.  

Imagine maintaining pressure on the gas pedal in the car steadily so that, starting from a stop at zero time, after $1$ second your velocity is $10$ mph, after $2$ seconds it is $20$ mph, after $3$ seconds, $30$ mph. If we continue at the same rate of acceleration, we'll go from $0$ to $60$ mph in $6$ seconds, which is quite a respectable time.

This example has constant acceleration.  Here, we say that $v$ is a constant function of time, and write 
\[ v = at \]
where $a$ is the acceleration.  It is the slope of the line.

\begin{center} \includegraphics [scale=0.35] {velocity_time_4.png} \end{center}
What about the distance?  It turns out that the distance is once again the area under the curve.

If $a$ is non-zero and constant, then $v$ changes at a constant rate.  Starting from $0$, the final velocity will be $v = at$, but the distance traveled is no longer the product 
\[ s = v \times t \stackrel{?}{=}  \]

because this $v$ is the final velocity and that is not the correct $v$ to use.  For variable velocity, the distance traveled is the \emph{average} velocity times the time.  For smooth (constant) acceleration from zero to $v$, the average velocity is the average of the initial and final velocities:
\[ v_{\text{avg}} = \frac{1}{2} \ (v_i + v_f) = \frac{1}{2} \ v \]

So the correct equation is:
\[ s = v_{\text{avg}}\ t  = \frac{1}{2} \ v \cdot t \]
and since $v = at$
\[ s = \frac{1}{2} at^2 \]

In this case, if we plot velocity as a function of time, we obtain a straight line that extends diagonally up with respect to the $x$-axis.  The distance traveled is the area under the curve, below the line and above the $x$-axis.  

The shape whose area is needed is a triangle.  This also accounts for the factor of $1/2$.

$\bar{v}$ is the average velocity during the interval between $0$ and $t$.  Since the acceleration $a$ is a constant, this is simply the average of the velocity at the end $(v = at)$, and at the start $(v = 0)$.
\[ \bar{v} = \frac{1}{2} (v_f + v_0) = \frac{1}{2} (at + 0) = \frac{1}{2} at \]
So that's why there is a factor of one-half in the basic formula.
\[ x = \frac{1}{2}at^2 + v_0 t + x_0 \]

\subsection*{back to gravity}

Again, the equation for gravity is
\[ h = \frac{1}{2}g t^2 + v_0 t + h_0 \]

This says that if we throw an object like a rock or a cannonball straight up into the air, the height at time $t$ depends on the acceleration due to gravity $g$, the initial velocity $v_0$ in the upward direction, and the initial height $h_0$ relative to the earth's surface.

For gravity we have only two directions in one dimension, up and down.  The constant $g$ is negative while $v_0$ and $h$ are positive, because the force of gravity points down while we have defined up as positive.  We can rewrite this as
\[ h = -16t^2 + vt + h_0 \]

Here the units of velocity $v$ are feet per second, $h$ is in feet, and $g$ is $-32$ feet per second per second.  (If this were a physics book we'd use meters, and then $g$ would be $-10$ meters per second per second.

Suppose the initial height is $h_0 = 96$, which means that you are almost $100$ feet in the air, perhaps on a platform or the top of a cliff, and the initial velocity is upward at $16$ feet per second.  That's a little more than 20 miles per hour.

Clearly the rock will go up to some maximum height and then come down.  The question asks first, at what time is the height a maximum?  Recall the formula for the vertex.
\[ m = -\frac{b}{2a} = -\frac{v_0}{2(-16)} = \frac{16}{32} = \frac{1}{2} \]

The rock reaches its maximum height at $1/2$ second and that height is
\[ h = -16(1/2)^2 + 16(1/2) + 96 = 112 \]

My little brother throws 20 miles per hour.  Suppose you could talk a professional baseball player into climbing up there with you and throwing vertically up at 64 feet per second.  A fun fact is that $15$ mph is exactly $22$ feet per second.
\[ 64 \ \frac{22}{15} \approx 95 \]
about 95 miles per hour, a fastball.

Then
\[ m = - \frac{v_0}{2(-16)} = \frac{-64}{-32} = 2 \]
Two seconds to the maximum height.

If he throws the rock straight up, it will land on his head, or yours, in a stiff breeze.  Suppose the player throws it outward a bit, but still up with the same velocity, so when it comes down it just misses his nose.  

At what time will it reach the ground?  We have
\[ h = 0 = -16t^2 + 64t + 96 \]
and we need the roots of this quadratic equation because we are asking for $h = 0$ which is ground level.

Do you remember where we said that the roots don't change if we factor out a constant?
\[ h = 0 = (-16)(t^2 - 4t - 6) \]

We can't factor this one into integers.  The times where $h = 0$ are, with $m = 2$ and $c = -6$:
\[ m \pm \ \sqrt{m^2 - c} = 2 \pm \ \sqrt{4 + 6} = 2 \pm \ 3.16 = - 1.16, 5.16 \]

This says that the rock hits the ground after 5 seconds or so.  

The negative root also can have a physical meaning under some conditions.  The equation doesn't know about the baseball player or the ladder.  All it knows is that we started our clock at zero with the rock at $96$ feet above the ground and a vertical velocity of $16$ feet per second upward.  We won't worry about that.

If we are not allowed to use the simplification gained by factoring $a$ and also told to use the quadratic equation, you will write
\[ -16t^2 + 64t + 96 = 0 \]

so
\[ x = \frac{-b \pm \ \sqrt{b^2 - 4ac}}{2a} \]
\[ x = \frac{-64 \pm \ \sqrt{64^2 -4(-16)(96)}}{2(-16)} \]

I get $10240$ for the discriminant and about $3.16$ for the square root after dividing by $32$.  That matches the previous answer.

We can also say something about the velocity.  The equation for velocity is
\[ v = gt + v_0 \]

The velocity at any time $t$ is equal to the initial value plus the change due to gravity.  Again, $g$ is negative while $v_0$ is positive.  Depending on time $t$, $v$ may be either positive or negative.  We had a factor of $1/2$ before, that is $(1/2)g = -16$, but $g = -32$.

There is one moment when the velocity is zero, when the height is a maximum and the rock stops going up and comes back down.

In this problem
\[ v = 0 = -32t + 64 \]
Again, $2$ seconds to maximum height.

At $t = 5.16$ the velocity is about
\[ v = -32 (5.16) + 64 = 100 \]

\subsection*{maximizing distance}
King Charles requires you to shoot a cannonball as far as possible.  The distance depends on the angle at which the projectile is fired (we ignore air resistance).  

This classic problem is usually solved by trigonometry, but we are going to do it directly.  The velocity is some combination of horizontal and vertical components $v_x$ and $v_y$

Using the Pythagorean theorem we have that
\[ v^2 = v_x^2 + v_y^2 \]

Aside from this relationship, the two component velocities act independently.  If not for gravity the horizontal distance $d$ would be 
\[ d = v_x t \]

and the clock would just keep ticking, so $d$ would continue to increase, as it would if the cannon were fired in outer space.

However, the clock must stop, because for reasonable values of the velocity, gravity will bring the cannonball back down to earth before it leaves the atmosphere.  That equation is
\[ y = -16t^2 + v_y t + 0 \]
We want the time when $y = 0$ so
\[ 16t^2 = v_y t \]

Since there are only two terms in this quadratic and they both contain $t$, we can simplify it by dividing out
\[ 16t = v_y \]
\[ t = \frac{v_y}{16} \]

Dividing both sides by $t$ in this way, we "lose" a solution, but that solution is $t = 0$, which we don't care about.  Now plug that expression for $t$ back into the first equation.
\[ d = \frac{v_x v_y}{16} \]

We need to choose $v_x$ and $v_y$ to maximize $d$.  We can ignore the factor of $1/16$, since the value of $v_x$ or $v_y$ giving a maximum for $d$ will also maximize $16d$.

In the end, we want to maximize
\[ v_x v_y = v_x \sqrt{v^2 - v_x^2} \]
where $v$ is a constant for a given cannon and amount of gunpowder, etc.

Just to explore a little, scale the problem so that $v = 1$, and come up with a simpler name for the variable $v_x$, like $s$.  Then take a look in Desmos
\begin{center} \includegraphics [scale=0.5] {quad4.png} \end{center}

The curve doesn't look like a standard parabola (it actually is one, but rotated by an angle).  It seems that the maximum comes at roughly $s = 0.7$.  That's a hint, because we know that $1/\sqrt{2} \approx 0.7$.

Without getting into the gory details, there is no reason to favor $x$ over $y$.  We expect they should be equal.
\[ v^2 = v_x^2 + v_y^2 = 2 v_x^2 \]
\[ v_x = \frac{v}{\sqrt{2}} \]

The maximum comes when $v_x = v_y$, that is, for an angle of $45^{\circ}$.  I'm told this isn't true in real life, because of air resistance.

Putting this another way, to maximize this:
\[ v_x \sqrt{v^2 - v_x^2} \]
we maximize the square:
\[ v_x^2 (v^2 - v_x^2) \]

$v_x^2$ will be some fraction $f$ of $v^2$.
\[ fv^2 (v^2 - fv^2) = v^2 \ [ \ f(1-f) \ ] \]
$v$ is a non-zero constant, so this problem really amounts to maximizing $f (1-f)$.  The quadratic is 
\[ -f^2 + f = 0 \]
The maximum is $m = -b/2a = - (1)/(2)(-1) = 1/2$.  Plug it into Desmos and take a look.

So $v_x^2$ is one-half of $v^2$ and equal to $v_y^2$.  Hence $v_x = v_y$.

\subsection*{inclined plane}

Think about a ramp or inclined plane at some angle $\theta$ to the vertical.  Place a block of lead on the ramp and it slides down to the bottom.  We ignore friction.  

The vertical height is $h$, and the length of the ramp is $L$.  It always helps to draw a picture.
\begin{center} \includegraphics [scale=0.4] {incline.png} \end{center}

It is a remarkable fact that forces (and most everything) can be converted easily between different coordinate systems.  

The force along the ramp, making the block move, is the component of gravity acting in that direction.  Using similar right triangles, it is easy to show that 
\[ \frac{h}{L} = \frac{g'}{g} \]
This is called the cosine of the angle $\theta$ but we don't need trigonometry for this problem.  

The point is that a part of the force of gravity still acts along an inclined plane, multiplied by a factor less than one and equal to $h/L$.  So our equation of motion is
\[ s = \frac{1}{2} g' t^2 + v_0 t + s_0 \]
where $g' = gh/L$.

When does the block reach the bottom?  Then, $s = L$ and
\[ L = \frac{1}{2} g' t^2 = 16 h/L \cdot t^2 \]
\[ t = \frac{L}{4 \sqrt{h}} \]

The velocity at the bottom is 
\[ v = g't = g h/L \cdot  \frac{L}{4 \sqrt{h}} = 8 \sqrt{h} \]
Perhaps surprisingly, the answer is independent of $L$.

Different ramp lengths give the same final velocity, though not the same time.

If you know simple physics we can do this problem in another way, namely by conservation of energy.
\[ P.E. = K.E. \]
The potential energy at the top is equal to the kinetic energy at the bottom.
\[ mgh = \frac{1}{2} mv^2 \]
\[ v^2 = 2gh = 64 h \]
\[ v =  8 \sqrt{h} \]
which is what we had before.

\subsection*{line and a parabola}

Suppose we have generic parabola and a generic line in the plane, of just the sort we've been studying.

Here are the two most common possibilities.  The line crosses the parabola at two points or it doesn't cross at all
\begin{center} \includegraphics [scale=0.40] {para31.png} \end{center}

The other two possibilities both have a single point of intersection.  The tangent line at a point (left panel), and a vertical line (right panel).  In the latter case, the line has an undefined slope.

\begin{center} \includegraphics [scale=0.40] {para32.png} \end{center}

The line that just touches a curve at a point is called the tangent to the curve at that point, and this is a major topic of study in calculus.

Here's an example with actual numbers.

\begin{center} \includegraphics [scale=0.40] {quad6.png} \end{center}

The parabola is in red with its vertex at $x = -b/2a = -1$ and roots at $x = -3, 1$.  Three lines are shown, one crosses the curve at two points (blue), one does not cross at all (orange) and one is the tangent to the curve (green), touching only the point $(0,-3)$.

We will show how to find the equation of the green line easily.

As a reminder, suppose we have two lines
\[ y = k_1x + u \]
\[ y = k_2x + v \]

If we want to know where they meet, we are looking for some point $(x,y)$ that solves both.  Since the $x$'s and the $y$'s are equal we get
\[ k_1 x + u = y = k_2 x + v \]
and we just solve for $x$.

The problem is to do the same thing with a quadratic and a line.  

We have the line $y = kx + u$ and the parabola $y = ax^2 + bx + c$ and we set the two $y$'s equal to obtain
\[ kx + u = y = ax^2 + bx + c \]

Group terms
\[ ax^2 + (b - k)x + (c - u) = 0 \]

This is another quadratic and can be solved in the usual way, namely the roots are
\[ x = m \pm \ \sqrt{m^2 - (c-u)/a)} \]

The crucial observation is this:  the line that is \emph{tangent} to the curve at $(x,y)$ only touches at that single point.  That's what tangent means.

Therefore, we are interested in the case where the equation we wrote has only one solution.  That only happens if the discriminant is zero.

The square root is what gives us two solutions.  If the square root goes away there is only one solution and we have therefore
\[ x = m = -\frac{b - k}{2a} \]
\[ 2ax = - b + k \]
\[ k = 2ax + b \]

This is the slope of the tangent line to a parabola at the point $(x,y)$.  The slope is proportional to $x$, which makes sense.  The curve slopes more steeply upward, as $x$ gets bigger.

I can't help pointing out that at the vertex, the slope is $k = 0$.  So that means
\[ 2ax + b = 0 \]
\[ x = - \frac{b}{2a} \]
We have calculated $m$ again (and got the same answer).

The tangent to the curve $ax^2 + bx + c$ has the slope $2ax + b$.  You might guess that the $2$ in the slope comes from the power of $2$ in the parabola, and you'd be right.

\subsection*{note about $i$}
We said that if $4ac > b^2$ we'll have the square root of a negative number.  We called the term under the square root, $D$, the discriminant.

What people did was to invent a new kind of number $i$ with the property that $i^2 = -1$, so $\sqrt{-1} = i$.

Then any negative square root can be written as $\sqrt{-N} = \sqrt{N} \cdot \sqrt{-1} = Ni$.

We had this as an example:  
\[ x^2 + 2x + 2 = 0 \]
$a = 1$, $b = 2$ and $c = 2$.  Then
\[ b^2 - 4ac = 4 - 8 = -4 \]
The roots are
\[ x = -1 \pm \ \frac{\sqrt{-4}}{2} = -1 \pm i \]

If we write the equation out using the roots
\[ (x - s)(x - t) \]
\[ = (x - (- 1 + i))(x - (-1 - i)) \]
\[ = ((x + 1) - i))((x + 1) + i) \]
we have four terms:
\[ = (x + 1)^2 + i (x+1) - i (x+1) - i^2 \]
The two terms in the middle cancel and the last one is $+1$:
\[ = (x + 1)^2 + 1 = x^2 + 2x + 2 \]

A real discussion about what $i$ \emph{means} will have to wait.  But the point to notice is that when we have roots containing $i$ they are always symmetric, it is always
\[ (p + qi)(p - qi) = p^2 + q^2 \]
And when multiplying out the terms with $i$ cancel.

\subsection*{extra}

You may want (or be required) to know the fancy version of the quadratic formula:
\[ x = -\frac{b}{2a} \pm \ \frac{\sqrt{b^2 - 4ac}}{2a} \]

It is relatively easy convert between the two forms.  Starting with
\[ x = m \pm \ \sqrt{m^2 - c/a} \]

substitute for $m$
\[ x = - \frac{b}{2a} \pm \ \sqrt{(\frac{b}{2a})^2 - c/a} \]

The key is now to put the last term $-c/a$ over the same denominator as the first term under the square root, $(2a)^2 = 4a^2$.  To do that, multiply by $4a$ on top and bottom
\[ x = - \frac{b}{2a} \pm \ \sqrt{(\frac{b}{2a})^2 - \frac{4ac}{4a^2}} \]

We have $4a^2$ in the denominator which can come out from under the square root and simplify to
\[ x = - \frac{b}{2a} \pm \ \frac{\sqrt{b^2 - 4ac}}{2a}  \]

However, I really prefer
\[ x = m \pm \ \sqrt{m^2 - c/a} \]

especially since you must already know the formula for $m$, and the $c/a$ part is visible in
\[ x^2 + \frac{b}{a} x + \frac{c}{a} =  0 \]


\end{document}